\SourceLeaf{173}\bigskip
\Needspace{10\baselineskip}
\centerline{[3.]}

\smallskip

\centerline{\textls{Ideen.}}

\smallskip

Die genaueste Bestimmung der Deklinationen der Sterne ließe sich in der Nähe des Äquators mit guten Repetitionstheodoliten vornehmen, indem man ihre Azimutunterschiede gegenüber terrestrischen Objekten zur Zeit der größten östlichen und westlichen Digressionen beobachtete. \hfill\mbox{1813 April 21.}

\smallskip

In der Theorie der Parallellinien sind wir jetzt noch nicht weiter als Euklid. Das ist der beschämende Teil der Mathematik, der früher oder später eine ganz andere Gestalt annehmen muss. \hfill\mbox{1813 April 27.}

\smallskip

Ebenso wie eine Klasse von Wahrheiten der höheren Arithmetik, auf reelle Zahlen beschränkt, in enger Verbindung mit den Kreisfunktionen steht, werden sich die allgemeineren Wahrheiten der H.\,A., auf imaginäre Größen ausgedehnt, gleichsam in den lemniskatischen Funktionen spiegeln. Diese Ideen öffnen die Pforte zu einem äußerst reichhaltigen Feld der Analysis.

\smallskip

Das Kristallprisma an den Fernrohren vergrößert deren Länge etwa um $\frac{1}{3}$
 der Prismenseite, beeinträchtigt aber etwas die vollkommene Farblosigkeit.

\bigskip
\hrule
\bigskip

\centerline{BEMERKUNG.}

\smallskip

\textsc{Schumacher} hat sich während des Winters 1808 bis 1809 in Göttingen aufgehalten. Über \textsc{Gauss}’ Gespräche mit ihm hat er Aufzeichnungen in einem Heft mit dem Titel »GAUSSIANA« gemacht. Aus diesem Heft ist die Notiz [2.] entnommen, während die Notiz [3.] auf einem einzelnen Zettel verzeichnet ist.
\begin{flushright}\textsc{Stäckel}.\end{flushright}

\bigskip
\hrule
